\documentclass[11pt,english]{article}

% \usepackage{footnote}
\usepackage{float}
\usepackage{url}

% High quality fonts
\usepackage[T1]{fontenc}

% Underlining text
\usepackage[normalem]{ulem}

% Math options
\usepackage{amsmath}
\usepackage{amsthm}
\usepackage{nicefrac}

% Theorem naming
\newtheorem{lem}{Lemma}
\newtheorem{proposition}{Proposition}

% Special tikz drawings
\usepackage{tikz}
\def\checkmark{\tikz\fill[scale=0.4](0,.35) -- (.25,0) -- (1,.7) -- (.25,.15) -- cycle;}

% Caption options
\usepackage{caption}
\captionsetup{labelfont=bf, margin=1cm, tableposition=top, figurewithin=none, tablewithin=none}

\raggedbottom

% For tables
\usepackage{booktabs}
\usepackage{siunitx}
\usepackage{makecell}

% For graphs
\usepackage[all,cmtip]{xy}

\usepackage{xstring}
\usepackage{ifthen}
\usepackage{catchfile}

\newcommand{\scalarinput}[1]{$\input{../output/tables/scalars/#1}$} % Keep minus signs
\newcommand{\absinput}[1]{%
    \CatchFileDef{\filecontents}{../output/tables/scalars/#1}{}% Read the file contents into \filecontents
    \IfBeginWith{\filecontents}{-}{% Check if the first character is a minus sign
        \StrGobbleLeft{\filecontents}{1}[\filecontents]% Remove the first character
    }{}%
    $\filecontents$% Output the processed contents
}
\newcommand{\scalarpctinput}[1]{%
  \CatchFileDef{\filecontents}{../output/tables/scalars/#1}{}% Read the file contents into \filecontents
  \IfBeginWith{\filecontents}{0.}{% Check if the first character is a minus sign
      \StrGobbleLeft{\filecontents}{2}[\filecontents]% Remove the first character
  }{}%
  $\filecontents$% Output the processed contents
}
% For xmark
\usepackage{pifont}% http://ctan.org/pkg/pifont
\newcommand{\xmark}{\text{\ding{55}}}

% Columns
\usepackage{multicol}

% No figures below footnote
\usepackage[bottom]{footmisc}

% Adding notes to table
\newcommand{\tablenote}[1]  {\par\vspace{3pt}{\raggedright \scriptsize \emph{Notes: }#1\par}\vspace{2pt}}

% Colors
% These are my colors -- there are many like them, but these ones are mine.
\definecolor{blue}{RGB}{0,92,216}
\definecolor{red}{RGB}{213,40,0}
\definecolor{yellow}{RGB}{240,228,66}
\definecolor{green}{RGB}{0,158,115}
\definecolor{grey}{RGB}{180,180,180}
\definecolor{darkgrey}{RGB}{105,105,105}
\definecolor{white}{RGB}{255,255,255}
\definecolor{purple}{RGB}{199,124,255}
\definecolor{orange}{RGB}{197,90,17}

\newcommand{\grey}[1]{\textcolor{grey}{#1}}
\newcommand{\red}[1]{\textcolor{red}{#1}}
\newcommand{\green}[1]{\textcolor{green}{#1}}
\newcommand{\blue}[1]{\textcolor{blue}{#1}}
\newcommand{\purple}[1]{\textcolor{purple}{#1}}
\newcommand{\orange}[1]{\textcolor{orange}{#1}}

% Useful math commands
\newcommand{\argmax}[0]{\operatornamewithlimits{argmax}}
\newcommand{\argmin}[0]{\operatornamewithlimits{argmin}}
\newcommand{\E}{\mathbb{E}}
\newcommand{\R}{\mathbb{R}}
\renewcommand{\d}{\mbox{ d}}
\newcommand{\Var}{\mbox{Var}}

% Font -- seems to need to be after things
\usepackage[sc]{mathpazo}

% Bibliography
\usepackage[authordate, backend=biber, maxcitenames=2, mincitenames=1, maxbibnames=99, uniquelist=false, uniquename=false]{biblatex-chicago}
\makeatletter
\renewcommand{\@makefntext}[1]{%
  \setlength{\parindent}{0pt}%
  \makebox[1.5em][r]{\textsuperscript{\@thefnmark}}#1}
\makeatother

% Remove extraneous fields
\AtEveryBibitem{%
  \clearfield{doi}%
  \clearfield{issn}%
  \clearfield{isbn}%
  \clearfield{url}%
  \clearfield{urldate}%
  \clearfield{eprint}%
  \clearfield{note}%
}

% Subfiles -- needs to be near end
\usepackage{subfiles}

% For crossing things out
\usepackage{centernot}
\newcommand{\notimplies}{\centernot\implies}

% Math commands
\renewcommand{\lvert}[1]{\left.#1\right|}
\renewcommand{\rvert}[1]{\left|#1\right.}

% For separating out referee comments
\newenvironment{refcomment}
  {\clearpage\begin{quotation}\color{blue}\singlespacing\small}
{\end{quotation}}

\newenvironment{refcommentnoclear}
  {\begin{quotation}\color{blue}\singlespacing\small}
{\end{quotation}}

% For marking LLM-generated text (displays in red for review)
\newenvironment{llm}
  {\color{red}}
  {}

% For side by side graphs
\usepackage{subcaption}

% Coloring tables
\usepackage{colortbl}

% For inputting specific things about the counterfactuals
\newcommand{\passval}[1]{\scalarinput{#1_1}}
\newcommand{\passabs}[1]{\absinput{#1_1}}
\newcommand{\passse}[1]{(\scalarinput{#1_1_se})}
\newcommand{\passsetext}[1]{\scalarinput{#1_1_se}}

\newcommand{\nopassval}[1]{\scalarinput{#1_0}}
\newcommand{\nopassabs}[1]{\absinput{#1_0}}
\newcommand{\nopassse}[1]{(\scalarinput{#1_0_se})}
\newcommand{\nopasssetext}[1]{\scalarinput{#1_0_se}}

% Margins
\usepackage{geometry}
% \geometry{verbose,tmargin=1in,bmargin=1in,lmargin=1.25in,rmargin=1.25in}
\geometry{verbose,tmargin=1in,bmargin=1in,lmargin=1in,rmargin=1in}
% Float separation
\setlength{\textfloatsep}{12pt plus 2.0pt minus 2.0pt}

% Make section sizes smaller
\usepackage{titlesec}
\titleformat*{\section}{\large\bfseries}
\titleformat*{\subsection}{\normalsize\bfseries}
\titleformat*{\subsubsection}{\normalsize\bfseries}
\titlespacing\section{0pt}{8pt plus 4pt minus 2pt}{8pt plus 2pt minus 2pt}
\titlespacing\subsection{0pt}{8pt plus 4pt minus 2pt}{8pt plus 2pt minus 2pt}
\titlespacing\subsubsection{0pt}{8pt plus 4pt minus 2pt}{8pt plus 2pt minus 2pt}

% Hyphenation options
% \righthyphenmin62
% \lefthyphenmin62
\finalhyphendemerits=200000

% Amit's section styles
\renewcommand\thesection{\Roman{section}}
\renewcommand\thesubsection{\Roman{section}.\Alph{subsection}}
\renewcommand\thesubsubsection{\Roman{section}.\Alph{subsection}.\arabic{subsubsection}}

% Manual citation issues
\hyphenation{Muk-har-lya-mov}
\hyphenation{Shc-her-ba-kov}

% Spacing options
\usepackage{setspace}
% \setstretch{3.0}
\onehalfspacing

% Cleaner footnotes
\usepackage{fancyhdr}
\pagestyle{fancy}
\fancyhead{}
\fancyfoot{}
\fancyfoot[C]{\thepage}
\renewcommand{\headrulewidth}{0pt}

% Hyperref
\usepackage[unicode=true,bookmarks=true,bookmarksnumbered=false,bookmarksopen=false, breaklinks=true,pdfborder={0 0 0},pdfborderstyle={},backref=false,colorlinks=false]{hyperref}

\begin{document}

\renewcommand*{\thefootnote}{\fnsymbol{footnote}}

\author{

  \makebox[\textwidth][c]{%
    \begin{minipage}{0.33\textwidth}\centering
Mark Egan\thanks{Harvard University, Harvard Business School.
Email: megan@hbs.edu.} \\ \vspace{-1em} Harvard University and NBER

    \end{minipage}
    \hfill
    \begin{minipage}{0.33\textwidth}\centering
Gregor Matvos\thanks{Northwestern University, Kellogg School of Management.
Email: gregor.matvos@kellogg.northwestern.edu.} \\ \vspace{-1em} Northwestern University and NBER

    \end{minipage}
    \hfill
    \begin{minipage}{0.33\textwidth}\centering
Amit Seru\thanks{Stanford GSB.
Email: aseru@stanford.edu.} \\ \vspace{-1em} Stanford University and NBER

    \end{minipage}
}

  \vspace{2em}

  % Second row of authors
  \makebox[\textwidth][c]{%
    \begin{minipage}{0.33\textwidth}\centering
Lulu Wang\thanks{Northwestern University, Kellogg School of Management.
Email: lulu.wang@kellogg.northwestern.edu.} \\ \vspace{-1em} Northwestern University

    \end{minipage}

    \begin{minipage}{0.33\textwidth}\centering
Vincent Yao\thanks{Georgia State University, J.
Mack Robinson College of Business.
Email: wyao2@gsu.edu.} \\ \vspace{-1em} Georgia State University

    \end{minipage}
} }

\title{Who Pays for Payments?}

\date{\today\thanks{We thank our discussants Matteo Benetton, Itamar Drechsler, German Gutierrez, Chuqing Jin, Scott Nelson, Marc Rysman, Zhu Wang, Carlo Wix, and seminar participants at the Bank of Italy, the Hong Kong Institute for Monetary and Financial Research (HKIMFR), Berkeley, the Carey Finance Conference, Emory, the Federal Reserve Board, Harvard, Imperial, Northwestern, NBER Household Finance, NBER Corporate Finance, the Stanford Institute for Theoretical Economics, Toulouse, and UNC Chapel Hill for helpful comments. Mark Egan is an unpaid member of the Fiserv Small Business Index Advisory Council. Vincent Yao is an unpaid member of the Fiserv Small Business Index Advisory Council and a paid consultant. Fiserv had no role in the design of the study, the analysis, the interpretation of the results, the writing of the paper, the decision to publish, and did not review or approve the content of the paper. }}

\maketitle

% \begin{center}
% \vspace*{-1cm}
% \textcolor{red}{PRELIMINARY AND INCOMPLETE.
% PLEASE DO NOT SHARE!}

% \end{center}
\pagenumbering{roman} \thispagestyle{empty} 
% \vspace{0.5cm}

\begin{abstract}
\begin{singlespace}

  % Longer version that clarifies the details behind these inputs
  % We use novel merchant-level data covering both card and cash transactions to quantify consumer cross-subsidies in U.S. consumer payments.
  We use novel data on the composition and cost of payments across U.S. merchants to quantify consumer redistribution in the payment system.
  Cards charge interchange fees to merchants to fund consumer rewards.
  When merchants raise prices for all consumers in response to these costs, users of low-cost payment methods (e.g., cash and debit) cross-subsidize high-reward credit card users who shop at the same merchant.
  This standard mechanism implicitly assumes that consumers using different payment methods shop at the same merchants and that merchants face similar fees. 
  We show instead that incidence depends on the joint distribution of payment choices across merchants.
  We document two key forces that shape redistribution.
  First, consumer sorting---where consumers who use different payment methods shop at different merchants---limits the exposure of cash and debit users to the effects of high interchange fees.
  Second, interchange fees vary across merchants; where users of different payment methods overlap, such as at large grocery stores, fees are lower due to sector discounts and private negotiations.
  We embed these forces in a sufficient-statistics framework that maps observed heterogeneity directly into redistribution.
  We estimate that interchange fees transfer approximately \$\absinput{combined_dollar_headline} billion every year from cash and debit users to credit card users.
  Consumer sorting and merchant fee heterogeneity reduce the magnitude of this regressive transfer by \absinput{pct_reduction}\%, but do not eliminate it.
  Finally, we show that both the Durbin Amendment and the rise of premium credit cards have been regressive, highlighting how policy and innovation can reshape the incidence of platform fees.

\end{singlespace}

\thispagestyle{empty} \setcounter{page}{0}\begin{flushleft}

Keywords: Interchange Fees, Redistribution, Durbin Amendment, Price Discrimination, Payment Systems

JEL Classification: E42, D14, L11, L81.\end{flushleft}

\end{abstract}
\clearpage{}

\pagenumbering{arabic} \renewcommand*{\thefootnote}{\arabic{footnote}}
\setcounter{footnote}{0}

% Include main sections
\subfile{introduction}
\subfile{institutional_setting}
\subfile{payment_costs}
\subfile{redistribution}
\subfile{results}
\subfile{conclusion}

% Bibliography
\clearpage{}
\begin{singlespace}
\begin{flushleft}
{\small
\bibliography{payment_methods,payments_fiserv_lulu}
% \bibliography{../payment_methods,../payments_fiserv_lulu}
}{\small\par}

\par\end{flushleft}
\end{singlespace}

% Include figures
\clearpage{}
\subfile{figures}

% Include tables
\clearpage{}
\subfile{tables}

% Include appendices
\appendix
\renewcommand{\thesubsection}{\Alph{section}.\arabic{subsection}}
\renewcommand{\thesubsubsection}{\Alph{section}.\arabic{subsection}.\arabic{subsubsection}}

\pagenumbering{arabic}
\renewcommand{\thepage}{A-\arabic{page}}
\renewcommand{\thetable}{A.\arabic{table}}
\renewcommand{\thefigure}{A.\arabic{figure}}
\setcounter{table}{0}  % Reset table counter
\setcounter{figure}{0}  % Reset figure counter

\clearpage{}
\subfile{appendix_theory}

\clearpage{}
\subfile{appendix_endogenous_choice}

\clearpage{}
\subfile{appendix_network_incidence}

\clearpage{}
\subfile{appendix_data}

\clearpage{}
\subfile{appendix_cash_model}

\clearpage{}
\subfile{appendix_structural}

\clearpage{}
\subfile{appendix_durbin}

\clearpage{}
\subfile{appendix_figures}

\clearpage{}
\subfile{appendix_tables}

\end{document}
